% Table: Balance tests (Chapter 4 appendix)
\begin{table}[htbp]
\centering
\caption{Balance tests: effect of A-level economics availability on pre-treatment outcomes}
\label{tab:balance}
\begin{tabular}{l cc}
\toprule
                                      & Coefficient    & SE             \\
Dependent variable                    & on availability&                \\
\midrule
Std.\ GCSE overall score             & -0.007             & (0.008)           \\
Std.\ GCSE maths score               & 0.025\sym{***}             & (0.005)          \\
Std.\ KS2 score (age 11)             & 0.001             & (0.004)           \\
\addlinespace
\% female                             & -0.004            & (0.0024)           \\
\% White British                      & -0.008\sym{**}             & (0.004)           \\
\% FSM eligible                       & 0.001             & (0.001)           \\
\% London/South East                  & 0.0004             & (0.0009)           \\
\addlinespace
School FE                             & \multicolumn{2}{c}{\checkmark}  \\
Cohort FE                             & \multicolumn{2}{c}{\checkmark}  \\
Observations                          & \multicolumn{2}{c}{1,335,629--1,409,039}          \\
\bottomrule
\end{tabular}
\begin{minipage}{\textwidth}
\smallskip
\footnotesize
\textit{Notes:} Each row reports the coefficient from a separate regression of the indicated pre-treatment variable on an indicator for A-level economics availability, with school and cohort fixed effects. Small coefficients support the assumption that within-school changes in availability are not driven by shifts in student composition; the two statistically significant rows (GCSE maths and the White British share) are discussed in Section~\ref{sec:framework}, and both variables are controlled directly in the estimating equations. Observations vary by row with the availability of the dependent variable. Standard errors clustered at the school level. Sample: GCSE cohorts 2002--2008, all schools. \sym{*} $p<0.10$, \sym{**} $p<0.05$, \sym{***} $p<0.01$.
\end{minipage}
\end{table}
